% Schéma du processus Scrum — Projet PFE (6 mois)
% Style inspiré du diagramme Scrum classique (bleu clair → bleu foncé)

\begin{figure}[h]
    \centering
    \resizebox{\textwidth}{!}{%
    \begin{tikzpicture}[
        phase/.style={
            circle,
            minimum size=2.4cm,
            line width=2.5pt,
            draw=scrumcyan,
            fill=scrumcyan!12
        },
        done/.style={
            circle,
            minimum size=2.4cm,
            line width=2.5pt,
            draw=scrumbnavy,
            fill=scrumbnavy
        },
        lbl/.style={
            text width=2.6cm,
            align=center,
            font=\small
        }
    ]

        % --- Phase de planification (bleu clair) ---
        \node[phase] (backlog) at (0, 0) {%
            \begin{tikzpicture}[scale=0.55, scrumcyan]
                \foreach \y in {0.35, 0, -0.35} {
                    \draw[line width=1.2pt, rounded corners=1pt] (-0.55,\y-0.08) rectangle (0.55,\y+0.08);
                }
            \end{tikzpicture}%
        };
        \node[lbl, below=0.15cm of backlog] {Backlog\\produit};

        \node[phase] (planning) at (3.6, 0) {%
            \begin{tikzpicture}[scale=0.55, scrumcyan]
                \foreach \y/\c in {0.35/false, 0/true, -0.35/false} {
                    \draw[line width=1.2pt, rounded corners=1pt] (-0.55,\y-0.08) rectangle (0.55,\y+0.08);
                    \if\c
                        \draw[line width=1.2pt] (-0.35,\y) -- (-0.18,\y-0.1) -- (-0.02,\y+0.12);
                    \fi
                }
            \end{tikzpicture}%
        };
        \node[lbl, below=0.15cm of planning] {Réunion de\\planification\\du sprint};

        \node[phase] (sprintbl) at (7.2, 0) {%
            \begin{tikzpicture}[scale=0.55, scrumcyan]
                \draw[line width=1.2pt, rounded corners=1pt] (-0.35,-0.45) rectangle (0.35,0.45);
                \draw[line width=1.2pt, rounded corners=1pt] (-0.15,-0.55) rectangle (0.55,0.35);
            \end{tikzpicture}%
        };
        \node[lbl, below=0.15cm of sprintbl] {Backlog\\du sprint};

        \draw[scrumcyan, line width=5pt, -{Triangle[fill=scrumcyan, length=8pt, width=10pt]}]
            (backlog.east) -- (planning.west);
        \draw[scrumcyan, line width=5pt, -{Triangle[fill=scrumcyan, length=8pt, width=10pt]}]
            (planning.east) -- (sprintbl.west);

        % --- Boucle Sprint (bleu foncé) — 6 mois au lieu de 1–4 semaines ---
        \begin{scope}[shift={(11.8, 0)}]
            \draw[scrumbnavy, line width=9pt, -{Triangle[fill=scrumbnavy, length=12pt, width=14pt]}]
                (30:3.1cm) arc (30:330:3.1cm);

            \node[align=center, text=scrumbnavy] at (0, -0.15) {%
                {\Large\bfseries SPRINT}\\[3pt]
                {\large 6 mois}\\[2pt]
                {\footnotesize (6 sprints)}
            };

            % Scrum quotidien (petite boucle)
            \draw[scrumbnavy, line width=4pt, -{Triangle[fill=scrumbnavy, length=6pt, width=8pt]}]
                (210:1.35cm) arc (210:510:1.35cm);
            \node[align=center, font=\scriptsize, text=scrumbnavy] at (0, 2.15)
                {\textbf{24H}\\Scrum quotidien};
        \end{scope}

        \draw[scrumcyan, line width=5pt, -{Triangle[fill=scrumcyan, length=8pt, width=10pt]}]
            (sprintbl.east) -- ++(0.9, 0)
            .. controls +(0.7, 0.55) and +(-1.0, 0.55) .. (8.7, 0.55);

        \draw[scrumbnavy, line width=7pt, -{Triangle[fill=scrumbnavy, length=10pt, width=12pt]}]
            (14.9, 0) -- (16.2, 0);

        % --- Travail terminé ---
        \node[done] (finished) at (17.8, 0) {%
            \begin{tikzpicture}[scale=0.55, white]
                \draw[line width=1.2pt, rounded corners=1pt] (-0.45,-0.55) rectangle (0.45,0.55);
                \draw[line width=1.4pt] (-0.22,-0.05) -- (0,0.18) -- (0.35,-0.35);
            \end{tikzpicture}%
        };
        \node[lbl, below=0.15cm of finished] {Travail\\terminé};

    \end{tikzpicture}%
    }
    \caption{Processus Scrum adapté au projet PFE sur une durée de 6 mois}
    \label{fig:scrum-process}
\end{figure}
